\begin{table}[htbp]
    \centering
    \caption{Fidelity of the synthetic datasets on the real evaluation set.
    Values are reported as mean $\pm$ standard deviation across three
    independently trained generator seeds.}
    \label{tab:rq1_fidelity}
    \small
    \setlength{\tabcolsep}{5pt}
    \begin{tabular}{lcccc}
        \toprule
        Generator
            & Wasserstein (scaled) $\downarrow$
            & \shortstack{Jensen--Shannon\\distance $\downarrow$}
            & Correlation $\downarrow$
            & Detector AUC $\rightarrow 0.5$ \\
        \midrule
        CTAB-GAN+
            & $\mathbf{0.1113 \pm 0.0083}$
            & $\mathbf{0.0063 \pm 0.0016}$
            & $\mathbf{0.0434 \pm 0.0003}$
            & $\mathbf{0.8968 \pm 0.0316}$ \\
        CTGAN
            & $0.2554 \pm 0.0314$
            & $0.1080 \pm 0.0126$
            & $0.1823 \pm 0.0049$
            & $0.9923 \pm 0.0018$ \\
        DP-CGANS
            & $0.1689 \pm 0.0239$
            & $0.0911 \pm 0.0134$
            & $0.1054 \pm 0.0025$
            & $0.9501 \pm 0.0151$ \\
        \bottomrule
    \end{tabular}

    \vspace{2mm}
    \begin{minipage}{0.96\linewidth}
        \footnotesize
        \textit{Notes:} Bold values indicate the best mean for each metric.
        Lower values are better for the three distance metrics. Detector AUC
        uses the same balanced probe protocol for every generator. The
        categorical metric is the mean square-root Jensen--Shannon distance
        with base-2 logarithms, not Jensen--Shannon divergence. Detector AUC
        is evaluated by proximity to 0.5, which represents chance-level
        real-versus-synthetic discrimination. DP-CGANS was evaluated in
        non-private mode.
    \end{minipage}
\end{table}
